% --- Anexo 1 ---
\chapter{\Large Guion de Entrevista Semiestructurada Preliminar a Docentes}
\label{anexo:entrevista_docentes}

\vspace{-1.5em}
\section*{Ficha técnica del instrumento}
\begin{itemize}
    \item \textbf{Instrumento:} 
        Guion de entrevista semiestructurada.
    \item \textbf{Población objetivo:}
         Plantel docente de la asignatura de Introducción a la Programación.
    \item \textbf{Propósito:} 
        Indagar sobre la situación actual de la asignatura de Introducción a la 
        Programación e identificar dificultades de aprendizaje, conceptos 
        complejos, errores comunes y percepciones docentes sobre la enseñanza 
        de la Programación Orientada a Objetos.
\end{itemize}

\blockquote{\textit{\textbf{Nota aclaratoria:} 
    Al tratarse de una entrevista semiestructurada, este guion sirvió como 
    guía base flexible. A lo largo del desarrollo de las sesiones con los 
    docentes, se formularon preguntas de profundización y repreguntas en 
    función de la dinámica de la conversación.}}

\vspace{-1em}
\section*{Cuestionario Base}

\subsection*{Bloque 1: Dificultades iniciales y aprendizaje de la programación}
\begin{enumerate}
    \item 
        ¿Cuáles son las principales dificultades que observa en los estudiantes 
        al iniciar programación?
    \item 
        En su experiencia, ¿qué conceptos suelen resultar más complejos de 
        comprender en las primeras etapas?
\end{enumerate}

\subsection*{Bloque 2: Dificultades específicas y errores frecuentes en POO}
\begin{enumerate}
    \setcounter{enumi}{2}
    \item 
        Específicamente en Programación Orientada a Objetos, ¿qué temas generan 
        mayor confusión?
    \item 
        ¿Qué errores son más frecuentes cuando los estudiantes comienzan a 
        aprender Programación Orientada a Objetos?
    \item 
        ¿A qué cree que se deben principalmente estas dificultades?
    \item 
        ¿Considera que el problema está más relacionado con la comprensión 
        conceptual o con la implementación en código?
\end{enumerate}

\subsection*{Bloque 3: Desafíos docentes, metodologías y evaluación}
\begin{enumerate}
    \setcounter{enumi}{6}
    \item 
        Como docente de la materia, ¿cuáles son los mayores desafíos que 
        enfrenta al enseñar programación básica y POO?
    \item 
        ¿Ha utilizado herramientas visuales o interactivas para enseñar 
        programación? ¿Qué resultados ha observado?
    \item   
        ¿Cómo determina si un estudiante realmente ha comprendido los conceptos 
        de PRogramación Orientada a Objetos más allá de solo 
        "hacer que el código funcione"?
    \item 
        ¿Tiene un estimado de cuántos estudiantes aprueban la materia, la 
        reprueban o la abandonan?
\end{enumerate}


%\chapter{Código fuente del sistema}
%\label{anexo:codigo}

